% GENERATED FILE. Edit canonical lessons, then run npm run generate:books.
% canonical-source-hash: fnv1a64:9ab83b94d2adb85c
% canonical-lessons: TE-C347-vindu, TE-C347-kanuka, TE-C347-puspaguccham, TE-C347-svagatincu, TE-C347-kalusuko

\chapter{Feasts and Presents}
\label{ch:te-feasts-and-presents}
\hlchaptermodality{347}

\begin{chapteropening}
\textbf{By the end of this chapter:} \emph{I can say a feast, a present, a bouquet, to greet warmly and to meet.}

Complete the last lesson of chapter 347: I can say a feast, a present, a bouquet, to greet warmly and to meet.
\end{chapteropening}

\section[\texorpdfstring{\te{విందు} (vindu)}{విందు (vindu)}]{\te{విందు} (vindu) — a feast; a banquet}

\label{lesson:TE-C347-vindu}



\begin{quote}
Before the new one: say the Telugu for a wedding pendant, then the Telugu for to invite.
\end{quote}

\subsection*{You'll want to know: \te{విందు}}
\textbf{\te{విందు}} — \emph{vindu} — \textquotedblleft{}a feast; a banquet\textquotedblright{}.

\begin{quote}
\te{పెళ్ళి} \te{తరువాత} \te{పెద్ద} \te{విందు} \te{జరిగింది}.
\end{quote}

\begin{quote}
\emph{peḷḷi taruvāta pedda vindu jarigindi.}
\end{quote}

\begin{quote}
After the wedding there was a big feast.
\end{quote}

\te{విందు} \te{ఇచ్చు}, to give a feast.

\begin{cousinweb}
\textbf{\te{విందు}} (\emph{vindu}) is native Dravidian: DEDR 5415 lists Telugu \emph{vindu}, \textquotedblleft{}a feast, a treat\textquotedblright{}, in one entry with \emph{vinta}, \textquotedblleft{}strange, new\textquotedblright{}, and Tamil \emph{viruntu}, \textquotedblleft{}a feast, a guest\textquotedblright{}.
\end{cousinweb}

\begin{grammarlens}[title={What you've built}]
One more word for telling what happened.
\end{grammarlens}

\subsection*{Your turn}
\begin{itemize}
\raggedright
  \item \emph{Say it:} \emph{vindu}
  \item \emph{Say it:} \emph{vindu}, once more
  \item \emph{Say it:} \emph{peḷḷi taruvāta pedda vindu jarigindi.} — After the wedding there was a big feast.
  \item \emph{Recall:} say the Telugu for a wedding pendant, then the Telugu for to invite, then say \emph{vindu} again
\end{itemize}

\begin{tcolorbox}[breakable,colback=teal!4,colframe=teal!35!black,title={Before you move on}]
What does \te{విందు} mean? (\textquotedblleft{}A feast; a banquet\textquotedblright{}.) Say it once more.
\end{tcolorbox}
\section[\texorpdfstring{\te{కానుక} (kānuka)}{కానుక (kānuka)}]{\te{కానుక} (kānuka) — a present; an offering}

\label{lesson:TE-C347-kanuka}



\begin{quote}
Before the new one: say the Telugu for to invite, then the Telugu for a feast.
\end{quote}

\subsection*{You'll want to know: \te{కానుక}}
\textbf{\te{కానుక}} — \emph{kānuka} — \textquotedblleft{}a present; an offering\textquotedblright{}.

\begin{quote}
\te{అతిథులు} \te{పెళ్ళికూతురికి} \te{కానుకలు} \te{ఇచ్చారు}.
\end{quote}

\begin{quote}
\emph{atithulu peḷḷikūturiki kānukalu iccāru.}
\end{quote}

\begin{quote}
The guests gave the bride presents.
\end{quote}

Plural: \te{కానుకలు}, \emph{kānukalu}.

Close to \te{బహుమతి}; \te{కానుక} is warmer, a gift to someone you honour.

\begin{cousinweb}
\textbf{\te{కానుక}} (\emph{kānuka}) is native Dravidian: DEDR 1443 lists Telugu \emph{kānuka}, \textquotedblleft{}a gift offered to a superior, a present\textquotedblright{}, under \emph{kanu}, \textquotedblleft{}to see\textquotedblright{}.
\end{cousinweb}

\begin{grammarlens}[title={What you've built}]
One more word for telling what happened.
\end{grammarlens}

\subsection*{Your turn}
\begin{itemize}
\raggedright
  \item \emph{Say it:} \emph{kānuka}
  \item \emph{Say it:} \emph{kānuka}, once more
  \item \emph{Say it:} \emph{atithulu peḷḷikūturiki kānukalu iccāru.} — The guests gave the bride presents.
  \item \emph{Recall:} say the Telugu for to invite, then the Telugu for a feast, then say \emph{kānuka} again
\end{itemize}

\begin{tcolorbox}[breakable,colback=teal!4,colframe=teal!35!black,title={Before you move on}]
What does \te{కానుక} mean? (\textquotedblleft{}A present; an offering\textquotedblright{}.) Say it once more.
\end{tcolorbox}
\section[\texorpdfstring{\te{పుష్పగుచ్ఛం} (puṣpagucchaṁ)}{పుష్పగుచ్ఛం (puṣpagucchaṁ)}]{\te{పుష్పగుచ్ఛం} (puṣpagucchaṁ) — a bouquet}

\label{lesson:TE-C347-puspaguccham}



\begin{quote}
Before the new one: say the Telugu for a feast, then the Telugu for a present.
\end{quote}

\subsection*{You'll want to know: \te{పుష్పగుచ్ఛం}}
\textbf{\te{పుష్పగుచ్ఛం}} — \emph{puṣpagucchaṁ} — \textquotedblleft{}a bouquet\textquotedblright{}.

\begin{quote}
\te{విద్యార్థులు} \te{అతిథికి} \te{పుష్పగుచ్ఛం} \te{ఇచ్చారు}.
\end{quote}

\begin{quote}
\emph{vidyārthulu atithiki puṣpagucchaṁ iccāru.}
\end{quote}

\begin{quote}
The students gave the guest a bouquet.
\end{quote}

\te{పుష్పం}, a flower, and \te{గుచ్ఛం}, a bunch.

\begin{cousinweb}
\textbf{\te{పుష్పగుచ్ఛం}} (\emph{puṣpagucchaṁ}) is borrowed from Sanskrit: \emph{puṣpa}, \textquotedblleft{}a flower\textquotedblright{} (CDIAL 8303), and \emph{guccha}, \textquotedblleft{}a bunch of flowers\textquotedblright{} (CDIAL 4172).
\end{cousinweb}

\begin{grammarlens}[title={What you've built}]
One more word for telling what happened.
\end{grammarlens}

\subsection*{Your turn}
\begin{itemize}
\raggedright
  \item \emph{Say it:} \emph{puṣpagucchaṁ}
  \item \emph{Say it:} \emph{puṣpagucchaṁ}, once more
  \item \emph{Say it:} \emph{vidyārthulu atithiki puṣpagucchaṁ iccāru.} — The students gave the guest a bouquet.
  \item \emph{Recall:} say the Telugu for a feast, then the Telugu for a present, then say \emph{puṣpagucchaṁ} again
\end{itemize}

\begin{tcolorbox}[breakable,colback=teal!4,colframe=teal!35!black,title={Before you move on}]
What does \te{పుష్పగుచ్ఛం} mean? (\textquotedblleft{}A bouquet\textquotedblright{}.) Say it once more.
\end{tcolorbox}
\section[\texorpdfstring{\te{స్వాగతించు} (svāgatiñcu)}{స్వాగతించు (svāgatiñcu)}]{\te{స్వాగతించు} (svāgatiñcu) — to greet warmly; to welcome}

\label{lesson:TE-C347-svagatincu}



\begin{quote}
Before the new one: say the Telugu for a present, then the Telugu for a bouquet.
\end{quote}

\subsection*{You'll want to know: \te{స్వాగతించు}}
\textbf{\te{స్వాగతించు}} — \emph{svāgatiñcu} — \textquotedblleft{}to greet warmly; to welcome\textquotedblright{}.

\begin{quote}
\te{గ్రామస్తులు} \te{అతిథులను} \te{సంతోషంగా} \te{స్వాగతించారు}.
\end{quote}

\begin{quote}
\emph{grāmastulu atithulanu santōṣaṅgā svāgatiñcāru.}
\end{quote}

\begin{quote}
The villagers welcomed the guests happily.
\end{quote}

Past (I …): \te{స్వాగతించాను}, \emph{svāgatiñcānu}.

The verb of \te{స్వాగతం}, a welcome.

\begin{cousinweb}
\textbf{\te{స్వాగతించు}} (\emph{svāgatiñcu}) is built on Sanskrit \emph{svāgata}, \textquotedblleft{}welcome\textquotedblright{}, literally \textquotedblleft{}well come\textquotedblright{} (CDIAL 13920), with the Telugu verb ending \emph{-iñcu}.
\end{cousinweb}

\begin{grammarlens}[title={What you've built}]
One more word for telling what happened.
\end{grammarlens}

\subsection*{Your turn}
\begin{itemize}
\raggedright
  \item \emph{Say it:} \emph{svāgatiñcu}
  \item \emph{Say it:} \emph{svāgatiñcu}, once more
  \item \emph{Say it:} \emph{grāmastulu atithulanu santōṣaṅgā svāgatiñcāru.} — The villagers welcomed the guests happily.
  \item \emph{Recall:} say the Telugu for a present, then the Telugu for a bouquet, then say \emph{svāgatiñcu} again
\end{itemize}

\begin{tcolorbox}[breakable,colback=teal!4,colframe=teal!35!black,title={Before you move on}]
What does \te{స్వాగతించు} mean? (\textquotedblleft{}To greet warmly; to welcome\textquotedblright{}.) Say it once more.
\end{tcolorbox}
\section[\texorpdfstring{\te{కలుసుకో} (kalusukō)}{కలుసుకో (kalusukō)}]{\te{కలుసుకో} (kalusukō) — to meet (each other)}

\label{lesson:TE-C347-kalusuko}



\begin{quote}
Before the new one: say the Telugu for a bouquet, then the Telugu for to greet warmly.
\end{quote}

\subsection*{You'll want to know: \te{కలుసుకో}}
\textbf{\te{కలుసుకో}} — \emph{kalusukō} — \textquotedblleft{}to meet (each other)\textquotedblright{}.

\begin{quote}
\te{పెళ్ళిలో} \te{నేను} \te{పాత} \te{స్నేహితులను} \te{కలుసుకున్నాను}.
\end{quote}

\begin{quote}
\emph{peḷḷilō nēnu pāta snēhitulanu kalusukunnānu.}
\end{quote}

\begin{quote}
At the wedding I met old friends.
\end{quote}

Past (I …): \te{కలుసుకున్నాను}, \emph{kalusukunnānu}.

Who you meet takes -\te{ను}: \te{అతన్ని} \te{కలుసుకున్నాను}. \te{మళ్ళీ} \te{కలుద్దాం} comes from the same verb.

\begin{cousinweb}
\textbf{\te{కలుసుకో}} (\emph{kalusukō}) is native Dravidian: it goes back to \emph{kalayu}, \textquotedblleft{}to join, unite, meet\textquotedblright{} (DEDR 1299), with the self-ending \emph{-kō}.
\end{cousinweb}

\begin{grammarlens}[title={What you've built}]
One more word for telling what happened.
\end{grammarlens}

\subsection*{Your turn}
\begin{itemize}
\raggedright
  \item \emph{Say it:} \emph{kalusukō}
  \item \emph{Say it:} \emph{kalusukō}, once more
  \item \emph{Say it:} \emph{peḷḷilō nēnu pāta snēhitulanu kalusukunnānu.} — At the wedding I met old friends.
  \item \emph{Recall:} say the Telugu for a bouquet, then the Telugu for to greet warmly, then say \emph{kalusukō} again
\end{itemize}

\begin{tcolorbox}[breakable,colback=teal!4,colframe=teal!35!black,title={Before you move on}]
What does \te{కలుసుకో} mean? (\textquotedblleft{}To meet (each other)\textquotedblright{}.) Say it once more.
\end{tcolorbox}
